\chapter{Index des méthodes}

\noindent\textit{Les méthodes du fascicule, regroupées par chapitre, pour retrouver
rapidement « comment faire » sans connaître le numéro du chapitre.}
\vspace{8pt}

\begin{multicols}{2}
\footnotesize

\textbf{\normalsize\color{onbo}Chapitre 1 --- Ressemblances et différences}
\begin{enumerate}[label=\arabic*.,itemsep=1pt,topsep=2pt,leftmargin=1.6em]
\item Reconnaître un caractère de l'espèce
\item Distinguer caractère héréditaire et caractère acquis
\end{enumerate}
\vspace{4pt}

\textbf{\normalsize\color{onbo}Chapitre 2 --- Les chromosomes}
\begin{enumerate}[label=\arabic*.,itemsep=1pt,topsep=2pt,leftmargin=1.6em]
\item Lire un caryotype
\item Identifier une trisomie sur un caryotype
\end{enumerate}
\vspace{4pt}

\textbf{\normalsize\color{onbo}Chapitre 3 --- Les gènes et la diversité humaine}
\begin{enumerate}[label=\arabic*.,itemsep=1pt,topsep=2pt,leftmargin=1.6em]
\item Interpréter un génotype
\item Construire un tableau de croisement
\end{enumerate}
\vspace{4pt}

\textbf{\normalsize\color{onbo}Chapitre 4 --- Transmission de l'information génétique}
\begin{enumerate}[label=\arabic*.,itemsep=1pt,topsep=2pt,leftmargin=1.6em]
\item Comparer mitose et fécondation
\item Relier un dérèglement cellulaire au cancer
\end{enumerate}
\vspace{4pt}

\textbf{\normalsize\color{onbo}Chapitre 5 --- Microorganismes et contamination}
\begin{enumerate}[label=\arabic*.,itemsep=1pt,topsep=2pt,leftmargin=1.6em]
\item Identifier une voie de contamination
\item Distinguer contamination et infection
\end{enumerate}
\vspace{4pt}

\textbf{\normalsize\color{onbo}Chapitre 6 --- Asepsie et antisepsie}
\begin{enumerate}[label=\arabic*.,itemsep=1pt,topsep=2pt,leftmargin=1.6em]
\item Distinguer une mesure d'asepsie d'une mesure d'antisepsie
\end{enumerate}
\vspace{4pt}

\textbf{\normalsize\color{onbo}Chapitre 7 --- La réponse immunitaire}
\begin{enumerate}[label=\arabic*.,itemsep=1pt,topsep=2pt,leftmargin=1.6em]
\item Reconnaître les signes d'une réaction inflammatoire
\item Distinguer réponse non spécifique et réponse spécifique
\end{enumerate}
\vspace{4pt}

\textbf{\normalsize\color{onbo}Chapitre 8 --- Le VIH/SIDA}
\begin{enumerate}[label=\arabic*.,itemsep=1pt,topsep=2pt,leftmargin=1.6em]
\item Situer une phase de l'infection à partir d'une description
\item Distinguer un mode de transmission réel d'une idée fausse
\item Proposer un message de prévention adapté
\end{enumerate}
\vspace{4pt}

\textbf{\normalsize\color{onbo}Chapitre 9 --- Les réactions allergiques}
\begin{enumerate}[label=\arabic*.,itemsep=1pt,topsep=2pt,leftmargin=1.6em]
\item Reconnaître une réaction allergique
\end{enumerate}
\vspace{4pt}

\textbf{\normalsize\color{onbo}Chapitre 10 --- Aider le système immunitaire}
\begin{enumerate}[label=\arabic*.,itemsep=1pt,topsep=2pt,leftmargin=1.6em]
\item Distinguer sérothérapie et vaccination
\end{enumerate}
\vspace{4pt}

\textbf{\normalsize\color{onbo}Chapitre 11 --- Circulation sanguine et accidents cardiovasculaires}
\begin{enumerate}[label=\arabic*.,itemsep=1pt,topsep=2pt,leftmargin=1.6em]
\item Distinguer infarctus et AVC
\item Relier un facteur de risque à un accident cardiovasculaire
\end{enumerate}
\vspace{4pt}

\textbf{\normalsize\color{onbo}Chapitre 12 --- La respiration et ses infections}
\begin{enumerate}[label=\arabic*.,itemsep=1pt,topsep=2pt,leftmargin=1.6em]
\item Retracer le trajet de l'air dans l'appareil respiratoire
\item Calculer une fréquence respiratoire
\item Interpréter une comparaison air inspiré / air expiré
\item Calculer un débit ventilatoire
\item Distinguer une infection respiratoire d'une maladie non infectieuse
\end{enumerate}
\vspace{4pt}

\textbf{\normalsize\color{onbo}Chapitre 13 --- Excrétion urinaire et insuffisance rénale}
\begin{enumerate}[label=\arabic*.,itemsep=1pt,topsep=2pt,leftmargin=1.6em]
\item Retracer le trajet de l'urine
\item Calculer un pourcentage de réabsorption rénale
\item Interpréter une analyse d'urine
\item Relier un facteur de risque à l'insuffisance rénale
\end{enumerate}
\vspace{4pt}

\textbf{\normalsize\color{onbo}Chapitre 14 --- Paludisme et fièvre Ebola}
\begin{enumerate}[label=\arabic*.,itemsep=1pt,topsep=2pt,leftmargin=1.6em]
\item Distinguer paludisme et fièvre Ebola
\end{enumerate}
\vspace{4pt}

\textbf{\normalsize\color{onbo}Chapitre 15 --- Les roches sédimentaires}
\begin{enumerate}[label=\arabic*.,itemsep=1pt,topsep=2pt,leftmargin=1.6em]
\item Retracer les étapes de formation d'une roche sédimentaire
\item Appliquer le principe de superposition
\item Utiliser un fossile pour reconstituer un milieu ancien
\end{enumerate}
\vspace{4pt}

\textbf{\normalsize\color{onbo}Chapitre 16 --- La protection des écosystèmes}
\begin{enumerate}[label=\arabic*.,itemsep=1pt,topsep=2pt,leftmargin=1.6em]
\item Reconnaître une relation d'interdépendance
\item Relier une activité humaine à une conséquence sur l'écosystème
\end{enumerate}
\vspace{4pt}

\end{multicols}
